\section{Síntesis de la propuesta}
Se propone comparar ANIA y L3C-DeepMFC en un banco de realimentación acústica en tiempo real, mediante verificación numérica, simulación y evaluación física. La investigación buscará determinar los compromisos entre estabilidad, conservación de voz, robustez y recursos computacionales, con trazabilidad de datos y configuraciones.

La infraestructura existente permitirá iniciar el piloto físico y desarrollar la reproducción neuronal, pero las conclusiones sobre desempeño se establecerán después de completar la campaña. El resultado de la tesis será una evaluación fundamentada de las condiciones de uso de ambos métodos y de su factibilidad en una plataforma embebida.
